\clearpage
\fancyhead[L]{\small ASTRONOMISCHES JAHRBUCH.}
\hypertarget{native-610-lower}{}
\begin{center}\small Göttingische gelehrte Anzeigen. Stück 81 und 82. Seiten 801..810. 21. Mai 1818.\\[3pt]\rule{.72\linewidth}{.3pt}\end{center}

\emph{Astronomisches Jahrbuch für das Jahr {\upshape 1820}, nebst einer Sammlung der neuesten Abhandlungen, Beobachtungen und Nachrichten aus dem Gebiet der astronomischen Wissenschaften. Mit Genehmigung der Königlichen Akademie der Wissenschaften berechnet und herausgegeben von} Dr. I. E. \textsc{Bode}, königlicher Astronom usw. Berlin. Beim Verfasser und in Kommission bei F. Dümmler. 256 Seiten im Oktavformat, nebst einer Kupfertafel.

\begingroup\brokenpenalty=10000\relax
Das Jahr 1820 zeichnet sich durch mehrere bemerkenswerte astronomische Erscheinungen aus. Außer der großen Sonnenfinsternis vom 7. September, die auch in unserer Gegend ringförmig sein wird, kommen mehrere Bedeckungen von Planeten durch den Mond vor. Davon werden eine des Mars am 28. Januar und zwei des Jupiter am 4. Juni und 18. Oktober in Berlin sichtbar sein. --- Die Zusätze beginnen mit den astronomischen Beobachtungen an der Berliner Sternwarte im Jahr 1816. Unter ihnen heben wir die Beobachtung der Sonnenfinsternis vom 19. November und die der Mondfinsternis vom 4. Dezember hervor. --- Beobachtungen in Halberstadt vom Hauptmann von \textsc{Wahl}. Die Breite von Halberstadt (Moritzplan) wurde durch Sonnenbeobachtungen mit einem 18-zölligen \textsc{Troughton}'schen Spiegelkreis zu $51^\circ\,54'\,5''9$ bestimmt. Herr \textsc{von Wahl} berechnet die Länge aus einer Sonnenfinsternis und einer Bedeckung des Aldebaran zu $1'\,17''6$ in Zeitmaß östlich der Seeberger Sternwarte. Die hier nach Herrn \textsc{von Wahl}s Rechnung für mehrere verglichene Orte bestimmten Konjunktionszeiten weichen um einige Sekunden von den Ergebnissen ab, die andere bei ihren Berechnungen erhielten. --- Sichtbare Helligkeitsänderungen Algols in den Jahren 1817---1819, berechnet von Prof. \textsc{Wurm}. --- Astronomische Anzeigen und Beobachtungen der Opposition des Jupiter 1816 vom Astronomen \textsc{Derflinger} in Kremsmünster. --- Astronomische Bemerkungen und Berechnung der Opposition der Vesta 1815 und der Sonnenfinsternis von 1820 von Dr. \textsc{Gerling} in Cassel. Dieser Artikel enthält auch eine Berichtigung eines übereilten Urteils über die atmosphärische Refraktion im Jahrbuch für 1819. --- Über das \textsc{Kepler}'sche Problem von Staatsrat und Ritter \textsc{von Schubert}. Dieser Aufsatz ist gewissermaßen eine ausführliche Entwicklung der \textsc{Laplace}'schen Lösung der Aufgabe durch eine unendliche Reihe, die nach den Sinuswerten der Vielfachen der mittleren Anomalie fortschreitet. Ihre Koeffizienten erscheinen ihrerseits als unendliche Reihen nach aufeinanderfolgenden Potenzen der Exzentrizität. In einem späteren Nachtrag hat der Verfasser die allerdings sehr mühsame Entwicklung bis zur 13. Potenz der Exzentrizität geführt. Insofern diese Entwicklung mathematisch interessant ist, muss man dem Verfasser dafür dankbar sein, auch wenn man den astronomischen Wert der Lösung durch eine solche unendliche Reihe im Vergleich mit sogenannten indirekten Lösungen ganz anders beurteilt als er und gerade das, was er zur Empfehlung der ersteren auf Kosten der anderen geltend \SourceLeaf{611}machen will, nur im umgekehrten Sinn für anwendbar hält. Denn gerade beim Gebrauch der Reihe kann man die Genauigkeit nur so weit steigern, wie man die Werte der numerischen Koeffizienten entwickelt hat, während die Genauigkeit der sogenannten indirekten Methode allein durch die Genauigkeit der verwendeten Sinustafeln bestimmt wird. Dass letztere, wenn sie zweckmäßig eingerichtet sind, bei einigermaßen beträchtlichen Exzentrizitäten unvergleichlich bequemer sind, ist ohnehin hinlänglich bekannt. --- Astronomische Beobachtungen an der Wiener Sternwarte 1816 von \textsc{Triesnecker} und \textsc{Bürg}. Beobachtungen von Verfinsterungen der Jupitermonde, Sternbedeckungen und einigen Planetenoppositionen sowie der Mondfinsternis vom 9. Juni und der Sonnenfinsternis vom 18. November. --- Über die Justierung des Meridianfernrohrs von Prof. \textsc{Littrow}. Dieser Aufsatz befasst sich mit der Aufgabe, die Lage des Meridianfernrohrs aus beobachteten Durchgängen zu bestimmen, wenn die Fehler eine endliche Größe haben. Eigentlich ist diese Aufgabe eher von mathematischem als von astronomischem Interesse, da in der Praxis nur solche Fehler vorkommen, die als unendlich klein anzusehen sind. Überdies hängt die Anwendung dieser Methode von der Kenntnis der wahren Kulminationszeiten ab, die man nur mithilfe eines anderen Instruments erhalten kann. Für eine eigenständige Justierung des Meridianfernrohrs dürfte man nur die Unterschiede der wahren Kulminationszeiten als bekannt ansehen und muss dann voraussetzen, dass die Achse des Instruments justiert oder die Neigung gegeben ist. Freilich eignet sich auch dieses Verfahren nur für Meridianfernrohre von geringerem Rang, und bei einem Instrument ersten Ranges ist eine Justierung, die von der Kenntnis der Rektaszensionen abhängt, unzulässig. --- Aus einem Brief desselben Astronomen wird ein Auszug aus einer Untersuchung wiedergegeben, die darauf abzielt, die Planetenmassen zu bestimmen, welche die Bewegung der Erde stören. Da die Untersuchung in der Zeitschrift für Astronomie ausführlich abgedruckt ist, erlauben wir uns einige Bemerkungen dazu. Der Vergleich mit der Zeitschrift zeigt, dass ein Druckfehler im vorliegenden Auszug die Unsicherheit (genauer den wahrscheinlichen Fehler) der Venusmasse auf das Zehnfache erhöht hat. Der Verfasser hat 189 Greenwicher Sonnenbeobachtungen mit \textsc{von Zach}s Sonnentafeln verglichen und durch Zusammenfassung in Gruppen 43 Bedingungsgleichungen entwickelt, die 11 unbekannte Größen enthalten Er beschränkt sich jedoch darauf, nach der Methode der kleinsten Quadrate nur für eine von ihnen, nämlich für die Korrekturen der Epoche, der Venusmasse, der Marsmasse und der Mondmasse, die Werte zu bestimmen. Wir sehen jedoch nicht recht ein, warum von den elliptischen Elementen gerade die Epoche allein und mit Vorrang vor der Länge des Perigäums und der Exzentrizität berücksichtigt wird. Auch die Aussage, dass die Epoche nicht leicht mit einiger Genauigkeit anzugeben sei, bedürfte wohl einer Korrektur. Ferner ist das Urteil (Zeitschrift, S. 285), dass sich die Nutation aus Sonnenbeobachtungen mit größerer Sicherheit bestimmen lasse als die Venusmasse, unhaltbar, zumindest wenn von Sonnenlängen die Rede ist, die aus beobachteten Rektaszensionen bestimmt sind. Unmittelbar beobachtet werden nämlich die Unterschiede der Rektaszension zwischen der Sonne und den Fundamentalfixsternen, und diese Unterschiede werden von der Nutation nur schwach beeinflusst. Die Sonnenlängen selbst enthalten schon schon die für die Fixsterne berechnete Nutation, und man muss daher immer nur nahezu genau dieselbe Nutation wiederfinden, die man der Rechnung zugrunde gelegt hatte. Was schließlich die Ergebnisse selbst betrifft, die Herr \textsc{Littrow} für die genannten vier Größen selbst ermittelt hat, so wäre eine Wiederholung der Rechnung sehr zu wünschen. In ihr sind nämlich mehrere Fehler erkennbar, bei denen sich schwer entscheiden lässt, ob es sich um Druckfehler, Schreibfehler oder Rechenfehler handelt. Aus der Natur der Methode der kleinsten Quadrate ergibt sich, dass die Koeffizienten von $\mu'''$ in der ersten und von $\mu$ in der dritten Gleichung auf S. 292 gleich sein müssen, da sie hier ganz verschieden sind. Auch die Koeffizienten von $\mu'''$ in der zweiten und von $\mu'$ in der dritten Gleichung sollen genau dieselben sein. Als Summe der Fehlerquadrate vor Anwendung \SourceLeaf{612}der gefundenen Korrektur ermittelte Herr \textsc{Littrow} 973, nach ihrer Anwendung 727. Uns scheint jedoch, dass diese Verringerung beträchtlich stärker hätte ausfallen müssen. Der Auffassung, es sei zweckmäßiger, die Fehler selbst mit ihren Vorzeichen zu addieren als ihre Quadrate, ist nicht zuzustimmen. Denn dass jene Summe verschwindet, belegt \emph{nur} die Wirksamkeit der Korrektur der Epoche und sagt gar nichts über die anderen Korrekturen aus. Wir wünschen, dass der Verfasser diese Bemerkungen als Zeichen der Aufmerksamkeit versteht, mit der wir seine wertvolle Untersuchung verfolgt haben, und dass er die Untersuchung seinem Versprechen gemäß bald in größerem Maßstab ausführt. --- Astronomische Beobachtungen aus dem Jahr 1816 an der Prager und an der Wilnaer Sternwarte. Zu den Prager Beobachtungen gehören die vollständige Beobachtung der Sonnenfinsternis vom 18. November und die Opposition der Ceres, zu den Wilnaer Beobachtungen die Opposition der Pallas. --- Besonders wertvoll sind die Beobachtungen von Professor \textsc{Bessel} an der Königsberger Sternwarte. Sie umfassen sechs Planetenoppositionen, fünf Sternbedeckungen, die Mondfinsternis vom 4. Dezember 1816, die Zenitdistanzen der Sonne zu beiden Sonnenwenden und die Rektaszension des Polarsterns. --- Bemerkungen anlässlich der großen Sonnenfinsternis am 19. November 1816. Da eine große Sonnenfinsternis eine der seltenen Gelegenheiten ist, bei denen die Wetterverhältnisse an einer großen Zahl von Orten öffentlich bekannt werden, hatte der Herausgeber den glücklichen Gedanken, alle Orte zusammenzustellen, an denen ein klarer Himmel die vollständige Beobachtung erlaubte, ferner jene, an denen wechselhaftes Wetter nur zeitweise etwas von der Finsternis zu beobachten erlaubte, und schließlich jene, an denen ein ganz bedeckter Himmel gar nichts davon sichtbar werden ließ. Das Ergebnis zeigt, dass alle diese Orte ohne Zusammenhang und Regelmäßigkeit bunt durcheinanderliegen, und eignet sich daher sehr dazu, den Glauben an spezielle Wetterprophezeiungen etwas zu erschüttern. Der Verfasser fügt noch einige Bemerkungen zu der Erfahrung hinzu, dass es bei totalen Finsternissen gewöhnlich nicht so dunkel wird, wie man erwartet. Was er aber als Versuch anführt, den bei totalen Finsternissen zuweilen beobachteten Ring um den Mond zu erklären, ist uns nicht ganz klar geworden. --- Beobachtungen der Verfinsterungen der Jupitersatelliten und der Sternbedeckungen an der Greenwicher Sternwarte von 1811---1814 sowie die ebenfalls dort angestellten Beobachtungen des großen Kometen von 1811. --- Neue Elemente der Vestabahn von Professor \textsc{Gerling}. --- Astronomische Beobachtungen an der Göttinger Sternwarte von Prof. \textsc{Gauss}. Mondfinsternis vom 4. Dezember 1816, Polhöhe der neuen Sternwarte aus 158 Beobachtungen des Nordsterns und Beobachtungen der Sommersonnenwende von 1817. Diese Beobachtungen ergaben dasselbe Resultat, das die Astronomen seit mehreren Jahren in Verlegenheit versetzt. Die Polhöhe aus Zirkumpolarsternen fällt um mehrere Bogensekunden größer aus als die aus Sonnenbeobachtungen, obwohl die genannten Beobachtungen ohne das sonst am Objektiv angebrachte Gegengewicht vorgenommen wurden, das einige als mögliche Ursache dieses Unterschieds verdächtigten. Seitdem diese Beobachtungen, die die Unhaltbarkeit dieser Vermutung beweisen, durch die Zeitschrift für Astronomie bekannt gemacht worden sind, hat Herr Prof. \textsc{Bohnenberger} eine neue Hypothese zur Erklärung des rätselhaften Phänomens aufgestellt. Daraufhin wurde an dem \textsc{Reichenbach}-Kreis, mit dem die obigen Beobachtungen angestellt wurden, eine Änderung vorgenommen, deren Wirkung bei den Beobachtungen der letzten Wintersonnenwende für die Richtigkeit der Hypothese von \textsc{Bohnenberger} zu sprechen scheint. Ob diese Hypothese aber ausreicht, um das Phänomen \emph{ganz} zu erklären, oder ob man nicht vielmehr das Zusammenwirken \emph{mehrerer} Ursachen annehmen muss, wird an anderer Stelle untersucht werden. --- Beobachtung der Sonnenfinsternis vom 19. November 1816 in Glatz durch General \textsc{von Lindener}. --- Nachrichten über die Lebensumstände des verstorbenen \textsc{Triesnecker} von Prof. \textsc{Bürg}. --- Bestimmungen für den Polarstern von Baron \textsc{von Lindenau}. Die sehr verdienstvolle Auswertung von achthundert Beobachtungen des Polarsterns hat mehrere interessante Ergebnisse erbracht. Neu bestimmt wurden die \SourceLeaf{613}Aberrationskonstante zu $20''44861$ (wahrscheinlicher Fehler nach der Wahrscheinlichkeitstheorie von \textsc{Gauss-Laplace} $\pm0''032$), die Nutationskonstante zu $8''97707$ (wahrscheinlicher Fehler $\pm0''04421$) und die Parallaxe des Polarsterns zu $0''14444$ (wahrscheinlicher Fehler $\pm0''05568$). Aus der Nutationskonstante in Verbindung mit der Präzessionskonstante hat Herr \textsc{von Lindenau} außerdem versucht, mithilfe einiger Formeln der Mécanique Céleste und einer ihm von \textsc{Gauss} mitgeteilten Bedingungsgleichung die Abplattung der Erde und die Masse des Mondes abzuleiten. Für Erstere fand er einen Wert nahe $\tfrac{1}{316}$. In die verwendeten Formeln und die Rechnung haben sich jedoch mehrere Fehler eingeschlichen, deren Berichtigung durch Herrn \textsc{von Lindenau} selbst an geeigneter Stelle zu erwarten ist. Es ist bloßer Zufall, dass das obige Ergebnis der Wahrheit so nahekommt. Wenn alle Fehler sachgerecht korrigiert werden, ergibt sich für die Abplattung etwa $\tfrac{1}{203}$. Dabei ist jedoch zu beachten, dass die von \textsc{Gauss} angegebene Bedingungsgleichung (aus der in der hier abgedruckten Form der Koeffizient $\tfrac{3}{5}$ weggelassen werden muss) insofern hypothetisch ist, als sie die Ähnlichkeit der Schichten gleicher Dichte im Erdsphäroid voraussetzt. --- Im Jahr 1817 an der Seeberger Sternwarte vom Adjunkten \textsc{Enke} beobachtete Oppositionen des Jupiter und des Uranus. Die letztere ist besonders deshalb wichtig, weil Uranus noch nie so nahe an seinem Knoten beobachtet worden ist. --- Über die Verbesserung einer bereits ungefähr bekannten Kometenbahn, von Dr. \textsc{Olbers}. Hier wird ein unangemessenes Urteil berichtigt, das ein Astronom über die von \textsc{Olbers} für den genannten Zweck vorgeschlagene Methode gefällt hatte. Zugleich wird gezeigt, wie sich diese Methode durch Hinzufügen einer vierten Hypothese zur Bestimmung der elliptischen Bahn einrichten lässt. Dieses brauchbare Verfahren verdient es, neben anderen empfohlen zu werden. Der Verfasser bemerkt im Übrigen zu Recht, dass es eine nicht zu rechtfertigende Verschwendung von Kraft und Zeit wäre, wenn man bei \emph{allen} Kometen die Bahn als Ellipse zu berechnen unternähme. Dieses Urteil darf jedoch niemanden, der dazu in der Lage ist, davon abschrecken, die Bahn eines Kometen, dessen beobachtete Bewegung eindeutig von der Parabel abweicht, mit aller Sorgfalt zu bestimmen. Auch wird der hochverdiente Verfasser es gewiss nicht missbilligen, wenn angehende Astronomen, um ihre Fähigkeiten für eine Aufgabe zu üben, die doch nicht jedermanns Sache ist, sich vielleicht auch an einem Kometen versuchen, bei dem das Ergebnis am Ende nur eine Abweichung von der Parabel zeigt, die innerhalb recht unsicherer Grenzen liegt. --- Neue Berechnung der Säkularänderungen der Elemente der Erdbahn, von Prof. \textsc{Nicolai}. Diese sehr verdienstvolle Arbeit wurde durch die Vermutung veranlasst, dass der Unterschied zwischen der Bestimmung der Venusmasse, die Herr \textsc{von Lindenau} mithilfe der Theorie des Merkur gefunden hat, und derjenigen, die aus der Abnahme der Schiefe der Ekliptik nach der Theorie von \textsc{Laplace} gefolgert wurde, darauf beruhen könnte, dass Letzterer nur eine angenäherte Bestimmung der Säkularänderungen gibt. Der Verfasser hat daher die Säkularänderungen sämtlicher Elemente der Erdbahn nach der neuen (noch nicht öffentlich bekannt gemachten) Methode von \textsc{Gauss}, die vollkommen streng ist, mit größter Sorgfalt neu berechnet. Das Ergebnis hat jene Vermutung jedoch nicht bestätigt, da die neuen Resultate von den älteren doch nicht wesentlich abweichen. --- Beobachtung der Oppositionen von Vesta, Pallas, Jupiter und Uranus im Jahr 1817 an der Mannheimer Sternwarte durch denselben Astronomen. --- Geometrischer Lauf der Vesta 1817---1818, berechnet von Prof. \textsc{Gerling} in Marburg, und der Pallas, berechnet von Dr. \textsc{Tittel} in Göttingen. --- Ein Auszug aus einem Schreiben von Dr. \textsc{Struve} in Dorpat berichtet unter anderem über eine noch vorzunehmende trigonometrische Vermessung Livlands. --- Preisverzeichnis astronomischer Instrumente im Institut des Herrn \textsc{von Utzschneider} in Benediktbeuern. --- Die vermischten astronomischen Beobachtungen, Bemerkungen und Nachrichten, die diesen Jahrgang abschließen, enthalten auch noch mehrere astronomische und geografische Bestimmungen, die wir aus Platzgründen hier nicht einzeln nennen können.
\par\clearpage\endgroup
